\begin{parts}

\part[\textbf{الف)}]
چون کارت‌ها فقط می‌توانند مقدارهای \(2\)، \(3\) و \(4\) داشته باشند، از حالت شروع فقط جمع‌های \(0,2,3,4,5\) قابل رسیدن‌اند. پس فضای حالت برابر است با:

\[
S=\{(0,C),(2,C),(3,C),(4,C),(5,C),\mathrm{Done}\}
\]

که \(C\) یعنی بازی هنوز در جریان است و \(\mathrm{Done}\) یعنی بازی تمام شده است.

\part[\textbf{ب)}]
اگر عامل کنش \lr{stop} را در حالت \((n,C)\) انتخاب کند، بازی تمام می‌شود و پاداش برابر \(n\) است:

\[
R((n,C),\mathrm{stop})=n
\]

اگر کنش \lr{draw} را انتخاب کند، پاداش آنی صفر است. اگر جمع جدید کمتر از \(6\) باشد، بازی ادامه پیدا می‌کند؛ وگرنه به حالت \(\mathrm{Done}\) می‌رود و امتیاز نهایی صفر می‌شود:

\[
R((n,C),\mathrm{draw})=0
\]

\[
R(\mathrm{Done},\cdot)=0
\]

\part[\textbf{ج)}]
برای مقدار بهینه داریم:

\[
V^*(s)=\max_a Q^*(s,a)
\]

و چون \(\gamma=1\)، محاسبات از حالت‌های آخر ساده‌تر انجام می‌شود.

برای حالت \((5,C)\):

\[
V^*(5,C)=\max\{5,0\}=5
\]

برای حالت \((4,C)\):

\[
V^*(4,C)=\max\{4,0\}=4
\]

برای حالت \((3,C)\):

\[
Q^*((3,C),\mathrm{draw})=\frac{1}{3}V^*(5,C)+\frac{1}{3}(0)+\frac{1}{3}(0)=\frac{5}{3}
\]

\[
V^*(3,C)=\max\left\{3,\frac{5}{3}\right\}=3
\]

برای حالت \((2,C)\):

\[
Q^*((2,C),\mathrm{draw})
=
\frac{1}{3}V^*(4,C)+\frac{1}{3}V^*(5,C)+\frac{1}{3}(0)
=
\frac{4+5}{3}=3
\]

\[
V^*(2,C)=\max\{2,3\}=3
\]

برای حالت \((0,C)\):

\[
Q^*((0,C),\mathrm{draw})
=
\frac{1}{3}V^*(2,C)+\frac{1}{3}V^*(3,C)+\frac{1}{3}V^*(4,C)
=
\frac{3+3+4}{3}
=
\frac{10}{3}
\]

\[
V^*(0,C)=\frac{10}{3}
\]

پس مقدارهای نهایی و سیاست بهینه به صورت زیر است:

\[
V^*(0,C)=\frac{10}{3},\qquad \pi^*(0,C)=\mathrm{draw}
\]

\[
V^*(2,C)=3,\qquad \pi^*(2,C)=\mathrm{draw}
\]

\[
V^*(3,C)=3,\qquad \pi^*(3,C)=\mathrm{stop}
\]

\[
V^*(4,C)=4,\qquad \pi^*(4,C)=\mathrm{stop}
\]

\[
V^*(5,C)=5,\qquad \pi^*(5,C)=\mathrm{stop}
\]

\[
V^*(\mathrm{Done})=0
\]

\part[\textbf{د)}]
در این بخش فرض می‌کنیم در همه‌ی نمونه‌ها، هر بار که کنش \lr{draw} زده شده، کارت \(2\) آمده است. پس از دید داده‌ها، \lr{draw} به شکل زیر عمل می‌کند:

\[
(n,C)\xrightarrow{\mathrm{draw}}(n+2,C)
\]

البته اگر \(n+2\ge 6\)، بازی تمام می‌شود. با تعداد خیلی زیاد نمونه، مقدارهای \(Q\) به مقدارهای زیر همگرا می‌شوند:

\[
Q((0,C),\mathrm{stop})=0,\qquad Q((0,C),\mathrm{draw})=4
\]

\[
Q((2,C),\mathrm{stop})=2,\qquad Q((2,C),\mathrm{draw})=4
\]

\[
Q((3,C),\mathrm{stop})=3,\qquad Q((3,C),\mathrm{draw})=5
\]

\[
Q((4,C),\mathrm{stop})=4,\qquad Q((4,C),\mathrm{draw})=0
\]

\[
Q((5,C),\mathrm{stop})=5,\qquad Q((5,C),\mathrm{draw})=0
\]

این مقدارها لزوماً برابر \(Q^*\) واقعی نیستند، چون نمونه‌ها نماینده‌ی واقعی محیط نیستند و کارت‌های \(3\) و \(4\) اصلاً دیده نشده‌اند.

\part[\textbf{ه)}]
در حالت \((4,C)\) داریم:

\[
Q((4,C),\mathrm{stop})=4
\]

\[
Q((4,C),\mathrm{draw})=3
\]

پس کنش حریصانه \(\mathrm{stop}\) است. با \(\epsilon=0.4\) و دو کنش ممکن داریم:

\[
P(\mathrm{stop})=(1-\epsilon)+\frac{\epsilon}{2}
=
0.6+0.2=0.8
\]

\[
P(\mathrm{draw})=\frac{\epsilon}{2}=0.2
\]

مقدار مورد انتظار \(Q\) بر اساس مقدارهای یادگرفته‌شده‌ی فعلی برابر است با:

\[
\mathbb{E}[Q]
=
0.8(4)+0.2(3)
=
3.8
\]

اگر مقدار واقعی \(Q^*\) را معیار بگیریم، چون در حالت \((4,C)\) با \lr{draw} حتماً می‌بازیم، داریم:

\[
\mathbb{E}[Q^*]
=
0.8(4)+0.2(0)
=
3.2
\]

\part[\textbf{و)}]
تعریف \lr{regret} هر گام به صورت زیر است:

\[
\mathrm{regret}_t=V^*(s_t)-Q^*(s_t,a_t)
\]

برای گام اول:

\[
s_1=(0,C),\qquad a_1=\mathrm{draw}
\]

\[
\mathrm{regret}_1
=
V^*(0,C)-Q^*((0,C),\mathrm{draw})
=
\frac{10}{3}-\frac{10}{3}
=
0
\]

برای گام دوم:

\[
s_2=(2,C),\qquad a_2=\mathrm{stop}
\]

\[
\mathrm{regret}_2
=
V^*(2,C)-Q^*((2,C),\mathrm{stop})
=
3-2
=
1
\]

برای گام سوم:

\[
s_3=(0,C),\qquad a_3=\mathrm{draw}
\]

\[
\mathrm{regret}_3
=
V^*(0,C)-Q^*((0,C),\mathrm{draw})
=
0
\]

برای گام چهارم:

\[
s_4=(4,C),\qquad a_4=\mathrm{draw}
\]

\[
\mathrm{regret}_4
=
V^*(4,C)-Q^*((4,C),\mathrm{draw})
=
4-0
=
4
\]

پس \lr{regret} تجمعی برابر است با:

\[
\mathrm{Regret}
=
0+1+0+4
=
5
\]

\[
\boxed{\mathrm{Regret}=5}
\]

\end{parts}